\documentclass[11pt,a4paper]{article}
\usepackage[T1]{fontenc}
\usepackage{lmodern}
\usepackage[margin=26mm]{geometry}
\usepackage{amsmath,amssymb,amsthm,booktabs,array}
\usepackage{microtype}
\usepackage[hidelinks]{hyperref}
\hypersetup{
  pdftitle={Six further OEIS congruence proofs},
  pdfauthor={Dev Goyal},
  pdfsubject={Version 3: AI-generated proofs, integral constructions, and Lean verification}
}
\setlength{\parskip}{3pt}
\setlength{\emergencystretch}{2em}
\raggedbottom
\allowdisplaybreaks[1]
\clubpenalty=10000
\widowpenalty=10000
\newtheorem{theorem}{Theorem}[section]
\newtheorem{lemma}[theorem]{Lemma}
\newtheorem{corollary}[theorem]{Corollary}
\theoremstyle{remark}
\newtheorem{remark}[theorem]{Remark}
\newcommand{\Z}{\mathbb Z}
\newcommand{\N}{\mathbb N}
\newcommand{\F}{\mathbb F}
\newcommand{\ps}[1]{#1[[x]]}
\newcommand{\rev}{^{\langle-1\rangle}}
\newcommand{\Plus}{\operatorname{Plus}}
\newcommand{\Minus}{\operatorname{Minus}}
\title{Six further OEIS congruence proofs\\
  \large Expanded proof package, version 3}
\author{Dev Goyal}
\date{24 October 2026}

\begin{document}
\maketitle

\begin{center}
\fbox{\begin{minipage}{0.90\linewidth}
\textbf{AI disclosure and review status.}
ChatGPT (OpenAI) generated the proofs, mathematical exposition, Lean
formalizations, verification programs, and AI reviews for this contribution
in response to Dev Goyal's requests. Dev Goyal is the sole human author and
initiated and requested the work. The OEIS sequences and conjectures are
credited to Paul D. Hanna. \textbf{Independent human review is pending.}
This paper does not claim OEIS acceptance, peer review, or worldwide priority.
\end{minipage}}
\end{center}

\begin{abstract}
We prove six coefficient assertions marked conjectural in the retrieved
OEIS entries A397588, A389536, A389537, A389542, A378580, and A378581.
The first has odd coefficients exactly at powers of two. Three further
series have sparse residue-two support modulo four, and two theta
substitution equations have residue-two support exactly at positive
squares. Every implicitly defined integer series is constructed and
identified by its literal source equation and a uniqueness theorem.
The methods are triangular coefficient recursion, genuine compositional
inversion, characteristic-two squaring, and controlled passage from
modulus two to modulus four. We additionally verify an already stated
divisibility-by-three property of A397588. All six target assertions have
unbounded Lean proofs; finite calculations independently check their
source correspondence. This paper supplements the two earlier papers
retained unchanged in the combined archive.
\end{abstract}

\section{Scope, notation, and source status}

All series are formal. We write $[x^n]A$ for a coefficient and $A\circ B$
for $A(B(x))$, with the outer series first. Such substitution is used only
when $B(0)=0$. A \emph{normalized} series is $x+O(x^2)$; a
\emph{constant-one} series is $1+O(x)$. The notation $O(x^r)$ means
divisibility by $x^r$, including over rings with zero divisors. Ordinary
powers and compositional powers are different operations. The symbol
$A\rev$ denotes compositional inversion, while $U^{-1}$ denotes the
multiplicative inverse of a unit. The set $\N$ includes zero.

The official pages were read on 24 October 2026
\cite{oeis588,oeis536,oeis537,oeis542,oeis580,oeis581}. Each explicitly
called the target in Table~\ref{tab:scope} a conjecture. Throughout the
table, $a_n=[x^n]A$, $k\in\N$, and
\[
 \Theta(x)=1+2\sum_{j\ge1}x^{j^2}.
\]
The first four series are normalized; the last two have constant
coefficient one. In each modulus-four row, all coefficients in the stated
range outside the displayed support are zero modulo four.

\clearpage
\begin{table}[ht]
\centering\small
\begin{tabular}{@{}lp{0.32\linewidth}p{0.44\linewidth}@{}}
\toprule
Entry & Source equation & Target assertion and range\\
\midrule
A397588 & $A=x+(xA^2)'$ &
  $a_n$ odd iff $n=2^k$, $n\ge1$.\\[3pt]
A389536 & $A(x^2+2x^3)=xA$ &
  $a_n\equiv2\pmod4$ iff $n=2^k+1$, $n>1$.\\[3pt]
A389537 & $A(A-x)=A^2$ &
  $a_n\equiv2\pmod4$ iff $n=3\cdot2^k$, $n>2$.\\[3pt]
A389542 & $A(A^2-x^2)=4A^3$ &
  $a_n\equiv2\pmod4$ iff $n=2^k+1$, $n>1$.\\[3pt]
A378580 & $A(x/A)=\Theta$ &
  $a_n\equiv2\pmod4$ iff $n$ is a square, $n>0$.\\[3pt]
A378581 & $A(xA)=\Theta$ &
  $a_n\equiv2\pmod4$ iff $n$ is a square, $n>0$.\\
\bottomrule
\end{tabular}
\caption{Six source conjectures addressed in this paper.}
\label{tab:scope}
\end{table}

The source label records the status of the retrieved page. Exact-identifier
searches with \emph{proof} or \emph{conjecture}, and searches for the defining
equations, did not identify an earlier proof of these targets. This bounded
check does not establish that a result was globally open or that this is its
first proof. The standard algebraic ingredients below are not claimed as
new discoveries.

One relevant prior result was checked separately. Bhave's proof, recorded
in the A196523 comments on 13 October 2026, treats
$F=x+xF^{\circ\ell}$ and gives coefficients one modulo $\ell-1$
\cite{bhave}. That displayed theorem does not directly state the theta
substitution results here; the inspected page contains no linked general
theta theorem. This comparison is limited to the inspected statement.

The source's divisibility-by-three assertion for A397588 is included as a
proved additional property, not as a seventh conjecture. The sign-run
observation in A389537 and the convergence question in A378580 are not
addressed. Supporting existence, transport, and coefficient lemmas are
proof infrastructure, not additional open problems.

\section{Formal-series tools}

\begin{lemma}[Coefficientwise construction]\label{lem:contraction}
Let $K$ be a commutative ring. Suppose an operator $\Phi$ preserves
normalized series and, for $r\ge2$, satisfies
\[
 G-H=O(x^r)\quad\Longrightarrow\quad
 \Phi(G)-\Phi(H)=O(x^{r+1}).
\]
Then $\Phi$ has exactly one normalized fixed point in $\ps K$.
More generally, a residual equation whose degree-$r$ condition determines
the next unknown coefficient with multiplier one has a unique solution
once its initial coefficients are fixed and its lower-degree conditions
hold identically.
\end{lemma}

\begin{proof}
Begin with $G_0=x$ and set $G_{m+1}=\Phi(G_m)$. Normalization gives
$G_1-G_0=O(x^2)$. Inductively,
$G_{m+1}-G_m=O(x^{m+2})$, so each coefficient eventually stabilizes.
Define $G$ using those stable coefficients. The displayed implication
ensures that each coefficient of $\Phi(G)$ equals the corresponding stable
coefficient, hence $G=\Phi(G)$. If two fixed points agree below degree
$r$, they agree through degree $r$; induction proves uniqueness.

For the residual version, set the next coefficient to the negative of the
known residual term. Its multiplier is one, so the construction uses only
ring operations. No earlier equation changes, and every equation is
eventually imposed. The same coefficient comparison proves uniqueness.
\end{proof}

\begin{lemma}[Units and reversion]\label{lem:inverse}
Every constant-one series has a unique constant-one multiplicative
inverse. Every normalized series has a unique normalized two-sided
compositional inverse. These inversions, formal differentiation, and
valid substitution commute with coefficient ring homomorphisms.
\end{lemma}

\begin{proof}
For $UV=1$ with $U(0)=1$, the coefficient of degree $n>0$ is $v_n$ plus
a finite expression in earlier coefficients of $V$, which determines
$v_n$ without division. For normalized $G$, the coefficient of degree
$n>1$ in $G(H)$ is $h_n$ plus an expression in earlier coefficients of
$H$, giving a unique normalized right inverse $H$.

To see the second inverse identity, construct a normalized right inverse
$J$ of $H$. Associativity gives
$G=G\circ(H\circ J)=(G\circ H)\circ J=J$, and hence $H\circ G=x$.
The same association proves uniqueness. Each coefficient of a valid
composition is a finite polynomial in coefficients. Differentiation
multiplies the degree-$(n+1)$ coefficient by $n+1$. Ring homomorphisms
therefore preserve these operations and equations; uniqueness identifies
the mapped inverses with the inverses of the mapped series.
\end{proof}

\begin{lemma}[Binary support]\label{lem:binary}
There is a unique normalized integral series $C$ satisfying $C=x+C^2$.
Its reduction modulo two is
\[
 P(x)=\sum_{k\ge0}x^{2^k}.
\]
More generally, for every integer $d\ge1$, the unique zero-constant
solution of $V+V^2=x^d$ in $\ps{\F_2}$ is
\[
 V=P(x^d)=\sum_{k\ge0}x^{d2^k}.
\]
\end{lemma}

\begin{proof}
The operator $G\mapsto x+G^2$ satisfies Lemma~\ref{lem:contraction},
since $G^2-H^2=(G-H)(G+H)$ and $G+H=O(x)$. In characteristic two,
every series $V$ satisfies $V(x)^2=V(x^2)$. This follows by truncating
above any fixed degree and applying the polynomial Frobenius identity.

Thus $V+V^2=x^d$ is a recursion in which the coefficient at an odd positive
index is fixed immediately, while the coefficient at an even index $n$
depends on the coefficient at $n/2<n$. The constant coefficient is fixed
to zero. This proves uniqueness. The proposed sum is well-defined
coefficientwise because its exponents tend to infinity; adding its square
cancels all terms except $x^d$. Set $d=1$ to identify the reduction of $C$.
\end{proof}

\begin{lemma}[Doubling and substitution]\label{lem:double}
For $U,V\in(\Z/4\Z)[[x]]$,
\[
 2U=2V\quad\Longleftrightarrow\quad
 U\equiv V\pmod2.
\]
If $B,D$ have zero constant coefficient and $B\equiv D\pmod2$, then
$2H(B)=2H(D)$ for every series $H$ over $\Z/4\Z$.
\end{lemma}

\begin{proof}
For residues $u,v$ modulo four, $2u=2v$ exactly when their reductions modulo
two agree. Apply this fact to each coefficient. Reduction modulo two
commutes with substitution, giving the second assertion. In particular,
this lemma never cancels two inside $\Z/4\Z$.
\end{proof}

These arguments are purely formal: there are no analytic convergence
assumptions. We now construct each source series before using its modular
properties.

\section{A397588: a differential equation and binary parity}

\begin{theorem}\label{thm:differential}
There is exactly one normalized $A\in\ps{\Z}$ satisfying
\begin{equation}\label{eq:differential}
 A=x+(xA^2)'.
\end{equation}
Its coefficients satisfy, for every $n\ge0$,
\[
 [x^n]A\text{ is odd}\quad\Longleftrightarrow\quad
 n=2^k\text{ for some }k\in\N.
\]
Also $A\equiv x\pmod3$, so $3\mid[x^n]A$ for every $n>1$.
\end{theorem}

\begin{proof}
Over any commutative ring, put $\Phi(G)=x+(xG^2)'$. If $G$ is normalized,
then $xG^2=O(x^3)$ and its derivative is $O(x^2)$, so $\Phi$ is normalized.
For normalized $G,H$ with $G-H=O(x^r)$,
\[
 x(G^2-H^2)=x(G-H)(G+H)=O(x^{r+2}).
\]
Formal differentiation lowers the order by at most one, including in
positive characteristic. Thus $\Phi(G)-\Phi(H)=O(x^{r+1})$.
Lemma~\ref{lem:contraction} constructs the unique solution over every
commutative ring, in particular over $\Z$. Its integer recursion is
\[
 a_0=0,\qquad a_1=1,\qquad
 a_n=(n+1)\sum_{j=1}^{n-1}a_ja_{n-j}\quad(n>1).
\]
This also exhibits integrality directly. Mapping coefficients preserves
the equation, and uniqueness identifies the reduction with the
corresponding canonical finite-ring solution.

In characteristic two the Leibniz rule gives $(A^2)'=2AA'=0$, hence
$(xA^2)'=A^2$. Equation~\eqref{eq:differential} becomes $A+A^2=x$.
Lemma~\ref{lem:binary} identifies $A$ modulo two with $P$, proving the
parity criterion. At $n=0$, both sides of that criterion are false.

In characteristic three, $x$ itself satisfies the equation because
$(x^3)'=3x^2=0$. Uniqueness gives $A=x$ in $\ps{\F_3}$, which proves the
additional divisibility assertion.
\end{proof}

The parity statement is the conjecture in A397588 \cite{oeis588}.
The divisibility by three was already stated there without a conjecture
label. It is included here with a proof and a stronger entire-series
identity, with no claim that it is an additional newly solved problem.

\section{A389536: a Mahler-type substitution modulo four}

\begin{theorem}\label{thm:mahler}
The source equation
\begin{equation}\label{eq:mahler}
 A(x^2+2x^3)=xA(x),\qquad A=x+O(x^2),
\end{equation}
has a unique solution in $\ps{\Z}$. With $C=x+C^2$ as in
Lemma~\ref{lem:binary}, it satisfies
\[
 A(x)\equiv x+2xC(x)\pmod4.
\]
Consequently, for every $n>1$,
\[
 [x^n]A\equiv
 \begin{cases}
  2\pmod4,&n=2^k+1\text{ for some }k\in\N,\\
  0\pmod4,&\text{otherwise}.
 \end{cases}
\]
\end{theorem}

\begin{proof}
Write $B=x^2+2x^3$. For $A=\sum_{n\ge1}a_nx^n$ with $a_1=1$, the
coefficient of $x^{n+1}$ in \eqref{eq:mahler}, for $n\ge2$, gives
\begin{equation}\label{eq:mahlerrec}
 a_n=
 \sum_{\substack{j\ge1\\0\le n+1-2j\le j}}
 a_j\binom{j}{\,n+1-2j\,}2^{n+1-2j}.
\end{equation}
Indeed, $B^j=x^{2j}(1+2x)^j$. Every contributing index obeys
$j\le(n+1)/2<n$, so this is a finite, division-free recursion in earlier
coefficients. The degree-two equation is already satisfied by $a_1=1$,
and the lower equations vanish. Thus the construction exists and is
unique over any commutative ring and commutes with reduction.

Work now modulo four. Denote the reduction of $C$ by the same letter,
so $C=x+C^2$, and put $D=x+2xC$. Its normalization follows from $C(0)=0$.
Modulo two, $B=x^2$, and Frobenius gives $C(B)=C(x^2)=C^2$.
Lemma~\ref{lem:double} therefore gives $2C(B)=2C^2$ modulo four. Compute
in $(\Z/4\Z)[[x]]$:
\begin{align*}
 D(B)
 &=B+2B\,C(B)\\
 &=x^2+2x^3+2x^2C(B)\\
 &=x^2+2x^3+2x^2C^2\\
 &=x^2+2x^2C=xD.
\end{align*}
The second line discards a multiple of four; the fourth uses $C^2=C-x$.
Thus $D$ satisfies the actual source equation in this ring. Uniqueness
identifies it with the reduction of the integer solution.

For $n>1$, its coefficient is $2[x^{n-1}]C$ modulo four.
Lemma~\ref{lem:binary} says this is two precisely when $n-1=2^k$;
otherwise it is zero.
\end{proof}

The source writes $A=A(x^2+2x^3)/x$ \cite{oeis536}.
Equation~\eqref{eq:mahler} is the equivalent division-free identity.
Multiplication by $x$ is injective in every formal power-series ring,
so this rewriting introduces no additional choice of solution.

\section{A389537: an alternating lacunary inverse}

\subsection{Construction and the original equation}

Over any commutative ring, define
\begin{equation}\label{eq:lacunaryR}
 R(x)=\sum_{k\ge0}(-1)^kx^{2^k}.
\end{equation}
This is a normalized formal series, and cancellation of successive terms
gives $R(x)+R(x^2)=x$. It is the unique normalized solution of this
equation: an odd degree above one has coefficient zero, and an even
degree is determined by minus the coefficient at half that degree.
Equivalently, $U\mapsto x-U(x^2)$ strictly increases the order of
differences between normalized inputs, so Lemma~\ref{lem:contraction}
applies. In Lean the latter construction is used.

Let $A=R\rev$ be the genuine normalized compositional inverse.
Substituting $A$ in $R+R(x^2)=x$ gives
\[
 x+R(A^2)=A.
\]
Apply $A$ to $R(A^2)=A-x$ to obtain the literal A389537 equation
\begin{equation}\label{eq:lacunarysource}
 A(A(x)-x)=A(x)^2.
\end{equation}
Conversely, let normalized $U$ satisfy \eqref{eq:lacunarysource},
and let $V=U\rev$. Applying $V$ on the outside gives
$U-x=V(U^2)$. Substituting $V$ for $x$ then gives
$x-V=V(x^2)$. The uniqueness of \eqref{eq:lacunaryR} implies $V=R$
and hence $U=A$. We have therefore constructed the unique normalized
integral solution of the original equation, not just an auxiliary
series with a related congruence. Coefficient maps preserve this
construction by Lemma~\ref{lem:inverse}.

\subsection{Parity, exact halving, and the modulus-four result}

\begin{theorem}\label{thm:lacunary}
For the preceding integer source series, $a_1=a_2=1$, and for every $n>2$,
\[
 a_n\equiv
 \begin{cases}
  2\pmod4,&n=3\cdot2^k\text{ for some }k\in\N,\\
  0\pmod4,&\text{otherwise}.
 \end{cases}
\]
\end{theorem}

\begin{proof}
Modulo two, Frobenius changes the inverse equation to $R+R^2=x$.
Thus $(x+x^2)\circ R=x$. The uniqueness of normalized compositional
inverses shows
\begin{equation}\label{eq:lacunaryparity}
 A\equiv x+x^2\pmod2.
\end{equation}
In particular, there is a well-defined integral series
\begin{equation}\label{eq:half}
 H=\frac{A-x-x^2}{2}\in\ps{\Z},\qquad H(0)=0.
\end{equation}
This division takes place coefficientwise in the integers, after the
divisibility proof \eqref{eq:lacunaryparity}; it is not division by two
inside $\Z/4\Z$. The initial values $a_1=a_2=1$ also follow directly
from $R=x-x^2+O(x^4)$ and $R(A)=x$.

Reduce \eqref{eq:half} modulo four, retaining the notation $H$, and set
$D=A-x=x^2+2H$. Both $D$ and $x^2$ have zero constant coefficient,
and they have the same reduction modulo two. Lemma~\ref{lem:double}
therefore gives $2H(D)=2H(x^2)$. Expanding the source equation yields
\begin{align*}
 A(D)&=D+D^2+2H(D)
       =x^2+2H+x^4+2H(x^2),\\
 A^2&=(x+x^2+2H)^2=x^2+2x^3+x^4.
\end{align*}
All equalities in this calculation are in $(\Z/4\Z)[[x]]$.
After subtraction,
\[
 2\bigl(H+H(x^2)\bigr)=2x^3.
\]
Applying the equivalence in Lemma~\ref{lem:double}, and writing
$\overline H$ for reduction modulo two, gives
\[
 \overline H+\overline H(x^2)=x^3,
 \qquad\text{equivalently}\qquad
 \overline H+\overline H^{\,2}=x^3.
\]
Since $\overline H(0)=0$, Lemma~\ref{lem:binary} identifies it uniquely:
\[
 \overline H=\sum_{k\ge0}x^{3\cdot2^k}.
\]
For $n>2$, equation \eqref{eq:half} gives $a_n\equiv2[x^n]H\pmod4$.
The displayed support proves both residue cases.
\end{proof}

This is the modulus-four conjecture in A389537 \cite{oeis537}.
The same entry's empirical observation about lengths of sign runs is a
separate problem and is not a consequence claimed here.

\section{A389542: an integral inverse that becomes an involution}

\subsection{An integral construction and exact source correspondence}

Construct a normalized series $L$ over any commutative ring from
\begin{equation}\label{eq:squareL}
 L=x+L^2-2xL(4x^3).
\end{equation}
The right side preserves normalization. If $G-H=O(x^r)$, then
$G^2-H^2=O(x^{r+1})$ for normalized inputs, while
$x(G-H)(4x^3)=O(x^{3r+1})$. Lemma~\ref{lem:contraction} proves
existence and uniqueness without division. More explicitly, for $n\ge2$,
\[
 \ell_n=\sum_{i=1}^{n-1}\ell_i\ell_{n-i}
 -\begin{cases}
  2\,4^{(n-1)/3}\ell_{(n-1)/3},&3\mid n-1,\\
  0,&3\nmid n-1,
 \end{cases}
 \qquad \ell_1=1.
\]
Set $R=x-2xL$. Equation~\eqref{eq:squareL} gives
\begin{align}
 R^2+R(4x^3)
 &=x^2+4x^2\bigl(-L+L^2+x-2xL(4x^3)\bigr)\notag\\
 &=x^2.\label{eq:squareR}
\end{align}
Let $A=R\rev$, which exists integrally by Lemma~\ref{lem:inverse}.
Composing \eqref{eq:squareR} on the right with $A$ gives
$x^2+R(4A^3)=A^2$. Applying $A$ to $R(4A^3)=A^2-x^2$ yields
the exact source equation
\begin{equation}\label{eq:squaresource}
 A\bigl(A(x)^2-x^2\bigr)=4A(x)^3.
\end{equation}

For completeness, this normalized integer source solution is unique.
If a normalized $U\in\ps{\Z}$ satisfies \eqref{eq:squaresource} and
$V=U\rev$, substitution of $V(t)$ gives
$U(t^2-V(t)^2)=4t^3$. Applying $V$ yields
\begin{equation}\label{eq:squareinverseunique}
 V(t)^2+V(4t^3)=t^2.
\end{equation}
If two normalized integer solutions of \eqref{eq:squareinverseunique}
first differ by $c x^r+O(x^{r+1})$, $r\ge2$, their squares first differ
by $2c x^{r+1}+O(x^{r+2})$. Their substituted difference has order at
least $3r\ge r+2$. Subtracting the equations forces $2c=0$ in $\Z$,
and hence $c=0$, a contradiction. Thus $V=R$ and $U=A$.
The uniqueness argument here is specifically over $\Z$; cancellation of
two is not asserted over rings with two-torsion.

\subsection{Reduction modulo four}

\begin{theorem}\label{thm:squareresidue}
The unique normalized integer solution of \eqref{eq:squaresource} satisfies,
for every $n>1$,
\[
 [x^n]A\equiv
 \begin{cases}
  2\pmod4,&n=2^k+1\text{ for some }k\in\N,\\
  0\pmod4,&\text{otherwise}.
 \end{cases}
\]
\end{theorem}

\begin{proof}
The construction of $L$ works uniquely over any commutative ring.
It therefore commutes with reduction, as do $R=x-2xL$ and its genuine
compositional inverse. Modulo two, equation \eqref{eq:squareL} becomes
$L+L^2=x$. Lemma~\ref{lem:binary} identifies its support with the powers
of two.

Now work modulo four. Since $R=x-2xL\equiv x\pmod2$, substitution gives
$L(R)\equiv L\pmod2$, and Lemma~\ref{lem:double} gives
$2L(R)=2L$. Consequently
\begin{align*}
 R(R)&=R-2R\,L(R)=R-2R\,L\\
     &=(x-2xL)(1-2L)=x-4xL+4xL^2=x.
\end{align*}
Thus $R$ is its own compositional inverse modulo four.
Uniqueness of compositional inverses identifies the reduction of $A$
with $R$:
\[
 A\equiv x-2xL\pmod4.
\]
For $n>1$, the coefficient is $-2\ell_{n-1}$ modulo four. This equals
two precisely when $\ell_{n-1}$ is odd, which holds exactly when
$n-1=2^k$; otherwise it is zero.
\end{proof}

The result is the single modulus-four conjecture in A389542
\cite{oeis542}. The entire involution argument uses multiplication by
two and reduction to characteristic two; it makes no cancellation of
two in $\Z/4\Z$.

\section{A378580 and A378581: theta substitution}

\subsection{A general integral construction}

The two source equations are
\begin{equation}\label{eq:thetaoriginal}
 A_-(x/A_-(x))=\Theta(x),\qquad
 A_+(xA_+(x))=\Theta(x),\qquad A_\pm(0)=1.
\end{equation}
We first solve them with an arbitrary constant-one target $T$ in place
of $\Theta$, over any commutative coefficient ring $K$.

\begin{lemma}\label{lem:normalroot}
For each $T\in\ps K$ with $T(0)=1$, there is exactly one normalized $G$
satisfying
\begin{equation}\label{eq:normalroot}
 G\circ G=GT.
\end{equation}
\end{lemma}

\begin{proof}
Let $E(G)=G\circ G-GT$. Its coefficients in degrees zero and one vanish
for every normalized $G$. Suppose $G-H=d x^r+O(x^{r+1})$, with $r\ge2$.
Decompose the difference of self-compositions as
\[
 G\circ G-H\circ H=(G-H)\circ G+(H\circ G-H\circ H).
\]
Because $G=x+O(x^2)$, the first summand is
$d x^r+O(x^{r+1})$. Write $H=x+x^2W$. The second summand is
$(G-H)+G^2W(G)-H^2W(H)$. The difference after $G-H$ has order at least
$r+1$: factor $G^2-H^2$, and use that substitution in any series
preserves an agreement modulo $x^r$. Thus the second summand also is
$d x^r+O(x^{r+1})$.

Finally, $(G-H)T=d x^r+O(x^{r+1})$, since $T(0)=1$. Therefore
\[
 E(G)-E(H)=d x^r+O(x^{r+1}).
\]
The degree-$r$ residual has multiplier $2-1=1$ on the next unknown
coefficient. Lemma~\ref{lem:contraction}'s triangular construction proves
existence and uniqueness with no division, over the stated ring.
\end{proof}

Write this unique $G$ as $G=xU$. Then $U(0)=1$ and
\[
 G\circ G=G\,U(G)=GT.
\]
Cancelling $G$ is valid even when $K$ has zero divisors: multiplication by
$x$ is injective coefficientwise, and $U$ is a unit. It follows that
$U(xU)=T$. Conversely, any constant-one solution of $U(xU)=T$ gives a
normalized solution $xU$ of \eqref{eq:normalroot}. Consequently there
exists exactly one such $U$; denote it by $\Plus(T)$.

Define
\[
 \Minus(T)=\Plus(T^{-1})^{-1},
\]
where both occurrences of inverse here are multiplicative unit inverses.
If $U=\Plus(T^{-1})$, then $U(xU)=T^{-1}$.
Substitution respects products, so $U^{-1}(xU)=T$. Since
$x/\Minus(T)=xU$, this is
\[
 \Minus(T)\bigl(x/\Minus(T)\bigr)=T.
\]
Conversely, if constant-one $V$ satisfies $V(x/V)=T$, then inversion of
that identity gives $V^{-1}(xV^{-1})=T^{-1}$. The uniqueness of
$\Plus(T^{-1})$ identifies $V^{-1}=U$ and hence $V=\Minus(T)$.
We have proved existence and uniqueness for both source equations
over every commutative ring, including $\Z$.

The quotient is literal: $V(x)\,(xV(x)^{-1})=x$. It is not interpreted as a
compositional inverse. Coefficient maps preserve the two source equations
and their unit inverses, so uniqueness also proves that $\Plus$ and
$\Minus$ commute with reduction of coefficients.

\subsection{The common square-support congruence}

\begin{theorem}\label{thm:theta}
Let $A_-=\Minus(\Theta)$ and $A_+=\Plus(\Theta)$ over $\Z$.
These are the unique constant-one integral solutions of
\eqref{eq:thetaoriginal}, and
\[
 A_-(x)\equiv A_+(x)\equiv\Theta(x)\pmod4.
\]
In particular, for either sign and every $n>0$,
\[
 [x^n]A_\pm\equiv
 \begin{cases}
  2\pmod4,&n=j^2\text{ for some }j\in\N,\\
  0\pmod4,&\text{otherwise}.
 \end{cases}
\]
\end{theorem}

\begin{proof}
Work over $\Z/4\Z$ and put $S=\sum_{j\ge1}x^{j^2}$, so $\Theta=1+2S$.
The series $B=x\Theta$ has zero constant coefficient and satisfies
$B\equiv x\pmod2$. Lemma~\ref{lem:double} gives $2S(B)=2S(x)$.
Therefore
\[
 \Theta(x\Theta(x))=1+2S(x\Theta(x))=1+2S(x)=\Theta(x).
\]
The uniqueness just proved identifies $\Plus(\Theta)=\Theta$ over this
ring. Also
\[
 \Theta^2=(1+2S)^2=1+4S+4S^2=1,
\]
so $\Theta^{-1}=\Theta$. It follows that
\[
 \Minus(\Theta)=\Plus(\Theta^{-1})^{-1}
 =\Plus(\Theta)^{-1}=\Theta^{-1}=\Theta.
\]
Compatibility with coefficient reduction identifies these series with
the reductions of the canonical integer solutions. Finally, the
coefficient of $\Theta$ at a positive index is two exactly when the
index is a square, and zero otherwise.
\end{proof}

This proves one original coefficient conjecture in each of A378580
and A378581 \cite{oeis580,oeis581}. Each equivalence is accompanied
by the complementary integer statement $4\mid a_n$ at nonsquares.
The constant term is one in both sequences; excluding $n=0$ is essential.
The modular argument only uses the form $1+2S$ and $S(0)=0$, so its
mechanism is more general than square support. No additional specializations
are counted as new problems in this paper.

\section{Formal coverage, corroboration, and reproduction}

The accompanying project uses Lean 4.19.0 and Mathlib revision
\begin{center}\small\ttfamily
c44e0c8ee63ca166450922a373c7409c5d26b00b.
\end{center}
The mathematical objects are actual Mathlib formal power series,
derivatives, substitutions, and coefficient homomorphisms. The modules
under \texttt{lean/Round3/} implement the arguments as follows:

\begin{center}\small
\begin{tabular}{@{}lp{0.58\linewidth}@{}}
\toprule
Module & Formal content\\
\midrule
\texttt{DifferentialParity} & Integral source solution, uniqueness,
  binary parity, and the additional modulus-three property.\\
\texttt{MahlerParity} & Literal source solution, uniqueness,
  comparison with $x+2xC$, and both residue clauses.\\
\texttt{LacunaryInverse} & Auxiliary construction, genuine inverse,
  both source-equation directions, exact integer halving, and support.\\
\texttt{SquareReversion} & Integral auxiliary solution, source
  identification and uniqueness, involution modulo four, and residues.\\
\texttt{ThetaSubstitution} & General constant-one source solutions,
  literal unit quotient, uniqueness, reduction, and both square criteria.\\
\texttt{ModFour} & Exact scalar and series doubling lemmas and
  justified coefficientwise integer halving.\\
\bottomrule
\end{tabular}
\end{center}

Separate typed contract files restate the literal equations, coefficient
domains, initial conditions, uniqueness claims, and unrestricted index
ranges. In particular, the quotient contract uses
\texttt{PowerSeries.invOfUnit} explicitly; reversion contracts use proved
two-sided compositional inverses. Thus the target sequences are not
replaced by unrelated series defined from the desired residue pattern.

The package's validation process rebuilds the first-party proof modules,
compiles the contracts, prints fully elaborated endpoint types, and audits
their transitive axiom dependencies. The permitted foundational axioms are
\texttt{propext}, \texttt{Classical.choice}, and \texttt{Quot.sound}.
The proof sources use no \texttt{sorry}, \texttt{admit},
\texttt{native\_decide}, or custom mathematical axiom. Finite identities
among residues are proved with kernel-checked \texttt{decide}; they are
ingredients in all-index series proofs, not substitutes for them.
The delivered validation report records the actual run and its endpoint
list. Endpoint counts measure proof coverage, not counts of open problems.

AI reviewers additionally examined the connection between the formal
definitions and the source statements. These semantic reviews are distinct
from kernel checking, and neither establishes historical novelty or human
endorsement. Independent human review remains pending.

\paragraph{Independent finite corroboration.}
Exact-integer programs construct coefficients from the original equations,
with separate auxiliary constructions where applicable. They match all
displayed source terms and check the following finite ranges:

\begin{center}\small
\begin{tabular}{@{}lp{0.69\linewidth}@{}}
\toprule
Entries & Recorded finite checks\\
\midrule
A397588 & Differential identity and parity through degree 500;
  divisibility by three from degree 2 through 500.\\
A389536 & Residue pattern and Catalan comparison through degree 2050;
  direct generic substitution through degree 81.\\
A389537 & Direct source recursion agrees with reversion of the
  alternating lacunary series through degree 80.\\
A389542 & Direct source recursion agrees with the auxiliary
  construction through degree 80; source and inverse equations
  checked through degree 81.\\
A378580--581 & Direct source recursions agree with the separate
  normalized-root construction through degree 64; both source
  equations and residue clauses checked.\\
\bottomrule
\end{tabular}
\end{center}

These finite ranges are evidence about transcription and source
identification. The coefficient conclusions proved in Lean are unbounded.
The machine-readable evidence is under \texttt{validation/}; individual
proof notes explain the exact checks and source prefixes.

\paragraph{Reproduction and retained work.}
From the extracted package, with the pinned Lean toolchain and Python
installed, run:
\begin{verbatim}
cd lean
lake exe cache get
lake build
cd ..
python3 scripts/check_formalization.py
\end{verbatim}
The project records dependency pins, explicit acceptance contracts,
source hashes, build logs, theorem types, and axiom-audit results.
The package-wide report is authoritative for the particular delivered run.

The combined archive retains the two earlier papers and their proof
modules unchanged. This paper adds the six assertions in
Table~\ref{tab:scope} and the separately identified additional property.
The user-supplied repository and its author details are credited
\cite{repo}; its earlier identifier or DOI is not represented as a DOI
for this expanded work. No submission to OEIS or publication acceptance
is claimed.

\begin{thebibliography}{10}
\raggedright
\bibitem{oeis588}
P. D. Hanna, OEIS A397588, 3 July 2026.
\url{https://oeis.org/A397588}. Accessed 24 October 2026.
\bibitem{oeis536}
P. D. Hanna, OEIS A389536, 28 October 2025.
\url{https://oeis.org/A389536}. Accessed 24 October 2026.
\bibitem{oeis537}
P. D. Hanna, OEIS A389537, 3 November 2025.
\url{https://oeis.org/A389537}. Accessed 24 October 2026.
\bibitem{oeis542}
P. D. Hanna, OEIS A389542, 1 November 2025.
\url{https://oeis.org/A389542}. Accessed 24 October 2026.
\bibitem{oeis580}
P. D. Hanna, OEIS A378580, 8 January 2025.
\url{https://oeis.org/A378580}. Accessed 24 October 2026.
\bibitem{oeis581}
P. D. Hanna, OEIS A378581, 8 January 2025.
\url{https://oeis.org/A378581}. Accessed 24 October 2026.
\bibitem{bhave}
A. Bhave, proof of the iterated-family congruence in the comments
of OEIS A196523, 13 October 2026.
\url{https://oeis.org/A196523}. Accessed 24 October 2026.
\bibitem{repo}
D. Goyal, \emph{oeis-a396843-proof}, user-supplied repository.
\url{https://github.com/dvgyl/oeis-a396843-proof}.
Original source revision and attribution are preserved in the package.
\bibitem{mathlib}
Mathlib contributors, formal proof library for Lean.
\url{https://github.com/leanprover-community/mathlib4}.
The exact dependency revision is pinned in
\texttt{lean/lake-manifest.json}.
\end{thebibliography}
\end{document}
